% ---- pg-346 （书页 334） ----

\chapter{主成分}
\phantomsection\label{chap:8}

\section{引言}
\phantomsection\label{sec:8.1}

主成分分析所关心的问题, 是通过一组变量的几个线性组合来解释这组变量的方差-协方差结构. 它的一般目的是: (1) 数据的压缩; (2) 数据的解释.

虽然要求 $p$ 个成分可以再现全系统的变异性, 但大部分变异性常常只用少数 $k$ 个主成分就可说明. 出现这种情况时, 这 $k$ 个主成分中所包含的信息和那 $p$ 个原变量所包含的(几乎)一样多. 于是这 $k$ 个主成分就可以用来取代那初始的 $p$ 个变量, 并且由对 $p$ 个变量的 $n$ 次测量值所组成的原数据集, 就压缩为对 $k$ 个主成分的 $n$ 次测量值所组成的数据集.

主成分分析常常揭示出一些先前不曾料想到的关系, 因而会对数据给出不同寻常的解释. 例 8.5 讨论的股票市场数据, 提供了这方面的一个好例子.

主成分分析更多的是一种达到目的的方法, 而非目的本身. 这是因为主成分分析频繁地用作许多大型调研的中间步骤, 例如, 可以将主成分分析插到多重回归(见第 7 章)或聚类分析(见第 12 章)中去; 此外还因为(尺度)主成分是第 9 章考虑的因子分析模型中协方差矩阵的一个“分解因子”.

\section{总体主成分}
\phantomsection\label{sec:8.2}

主成分, 在代数上是 $p$ 个随机变量 $X_1,X_2,\cdots,X_p$ 的一些特殊的线性组合, 而在几何上这些线性组合代表选取一个新坐标系, 它是以 $X_1,X_2,\cdots,X_p$ 为坐标轴的原坐标系旋转后得到的. 新坐标轴代表数据变异性最大的方向, 并且提供对协方差结构的一个较为简单但更精炼的刻画.

我们将会看到, 主成分只依赖于 $X_1,X_2,\cdots,X_p$ 的协方差矩阵 $\boldsymbol{\Sigma}$ (或相关矩阵 $\boldsymbol{\rho}$). 它们的进一步讨论不需要多元正态的假定. 另一方面, 对由多元正态总体导出的主成分可用常密度椭球来作有用的阐明. 进而, 当总体是多元正态时, 可由样本分量作出推断(见 8.5 节).

设随机向量 $\mathbf{X}'=[X_1,X_2,\cdots,X_p]$ 有协方差矩阵 $\boldsymbol{\Sigma}$, 其特征值 $\lambda_1\geqslant\lambda_2\geqslant\cdots\geqslant\lambda_p\geqslant 0$.

考虑线性组合
\phantomsection\label{eq:8-1}
\begin{equation}
\begin{aligned}
Y_1&=\mathbf{a}_1'\mathbf{X}=a_{11}X_1+a_{12}X_2+\cdots+a_{1p}X_p\\
Y_2&=\mathbf{a}_2'\mathbf{X}=a_{21}X_1+a_{22}X_2+\cdots+a_{2p}X_p\\
&\qquad\vdots\\
Y_p&=\mathbf{a}_p'\mathbf{X}=a_{p1}X_1+a_{p2}X_2+\cdots+a_{pp}X_p
\end{aligned}
\tag{8-1}
\end{equation}
则由\hyperref[eq:2-45]{式(2-45)}, 我们得到
